\chapter{Age-Related Macular Degeneration}
\index{Age-Related Macular Degeneration}
\label{sec:geriatric:Age-Related Macular Degeneration}

\section{Overview}
This condition affects many people worldwide and can range from mild to severe. Recognizing the condition early and managing it properly gives you the best chance for a good outcome. Understanding what causes this condition helps your doctor choose the right treatment for you. Learning about your condition and taking an active role in your care are key to managing it long term.

\begin{itemize}[noitemsep]
  \item Wet (exudative/neovascular) AMD is less common but more aggressive, driven by choroidal neovascularization. Treatment for dry AMD includes the AREDS2 formulation (vitamins C and E, lutein/zeaxanthin, zinc) for intermediate disease
  \item Wet AMD doctors typically treat it with intravitreal anti-VEGF injections.
\end{itemize}

\section{Causes}
This condition develops from a combination of genetic and environmental factors that disrupt normal body processes. Several factors can work together to trigger or make the condition worse. Knowing the specific cause helps your doctor choose the most effective treatment.

\section{Risk Factors}

\begin{itemize}[noitemsep]
  \item Genetic predisposition and family history
  \item Advancing age
  \item Lifestyle factors (diet, exercise, smoking, alcohol)
  \item Environmental exposures
  \item Underlying medical conditions
  \item Medication use
\end{itemize}

\section{Signs and Symptoms}

\begin{itemize}[noitemsep]
  \item Pain or discomfort in the affected area
  \item Difficulty performing daily activities
  \item Changes in how the affected area looks or works
  \item General symptoms may include fatigue, fever, or weight changes
  \item Symptoms may develop slowly or appear suddenly
\end{itemize}

\section{Progression and Stages}
This condition can progress through different stages over time. Early stages often involve mild symptoms that you may not notice right away. As the condition progresses, symptoms become more noticeable and may affect your daily activities. Advanced stages may involve more serious symptoms and complications.

\begin{itemize}[noitemsep]
  \item Age-related macular deterioration (AMD) is the leading cause of irreversible central vision loss in adults over 65. Dry (nonexudative) AMD accounts for 85--90\% of cases, characterized by drusen, retinal pigment epithelium shrinkage or wasting, and photoreceptor deterioration. outlook is guarded
  \item Low-vision rehabilitation is essential for all stages.
\end{itemize}

\section{Complications}

\begin{itemize}[noitemsep]
  \item The condition may worsen over time if not treated
  \item Increased risk of other infections or injuries
  \item Side effects from medications or other treatments
  \item Reduced quality of life
  \item Emotional and psychological effects
\end{itemize}

\section{Diagnosis}

\begin{itemize}[noitemsep]
  \item Your doctor will take a detailed medical history and perform a physical exam
  \item Lab tests to check how your organs are working
  \item Imaging tests (like X-rays or MRI) to look at the affected area
  \item Specialized tests depending on which body system is involved
  \item Referral to a specialist for further evaluation
\end{itemize}

\section{Prevention}

\begin{itemize}[noitemsep]
  \item Eat a balanced diet and exercise regularly
  \item Avoid known triggers and risk factors
  \item Go for regular check-ups so any problems are caught early
  \item Follow your doctor's screening recommendations
\end{itemize}

\section{Treatment Options}

\begin{itemize}[noitemsep]
  \item Medications to control symptoms and treat the underlying cause
  \item Lifestyle changes and self-care
  \item Physical therapy and rehabilitation
  \item Surgery when needed
  \item Regular follow-up care
\end{itemize}

\section{Living with It}
In older adults, AMD severely impacts quality of life---loss of reading, face recognition, and driving independence---and is linked to depression and falls.

\begin{itemize}[noitemsep]
  \item Follow your treatment plan as prescribed
  \item Keep up with regular appointments with your healthcare provider
  \item Adjust your daily habits as needed
  \item Reach out to family, friends, or support groups for help
  \item Keep learning about your condition and how to manage it
\end{itemize}

\section{Outlook}
Your outlook depends on the specific condition, how advanced it is when found, and how well it responds to treatment. Getting treated early generally leads to better results. Many people manage this condition effectively with the right treatment. Regular check-ups are important to monitor your condition and adjust treatment as needed.
